%% ============================================================
%% cap03.tex — Capítulo 3
%% Modelado UML del dominio académico-administrativo
%% ============================================================

\chapter{Modelado -UML- del dominio académico-administrativo}
\label{cap:uml}

\noindent
El modelado mediante el Lenguaje Unificado de Modelado (-UML-,
\textit{Unified Modeling Language}) constituye el puente entre el análisis
del dominio, desarrollado en el Capítulo~\ref{cap:requisitos}, y la
implementación, que se abordará en los Capítulos siguientes. Puesto que
el dominio del Estatuto General y del Estatuto Académico de la \ud\
presenta una estructura jerárquica y relacional de considerable complejidad,
el rigor en el modelado no es un requisito formal: es una condición
necesaria para que el código resultante sea comprensible, mantenible
y extensible.

Dos advertencias antes de comenzar. La primera: un diagrama -UML- no es
---o no debería ser--- una representación fotográfica del código que se
implementará. Es una herramienta de comunicación y de razonamiento; modela
aspectos del dominio que el equipo de desarrollo necesita discutir antes
de codificar. La segunda: \plantuml\ ---empleado a lo largo de este capítulo
como herramienta de generación de diagramas a partir de texto plano---
introduce limitaciones notacionales respecto del estándar -UML-
completo, pero su ventaja de ser versionable en sistemas de control de
código fuente la hace especialmente adecuada para el entorno pedagógico
de este libro.\footnote{%
  \plantuml\ implementa un subconjunto del estándar -UML- 2.5. Las limitaciones
  más relevantes son: ausencia de soporte nativo para algunos tipos de
  multiplicidad con restricciones, notación simplificada para interfaces
  y realización, y representación limitada de restricciones -OCL-. Para el
  nivel introductorio de este libro, estas limitaciones no generan
  menoscabo apreciable.}

%% ─────────────────────────────────────────────
\section{Introducción al lenguaje -UML- y \plantuml}
\label{sec:uml-intro}
%% ─────────────────────────────────────────────

El estándar -UML- 2.5 define catorce tipos de diagramas organizados en
dos grandes familias: diagramas estructurales y diagramas de comportamiento.
La Tabla~\ref{tab:diagramas-uml} describe los empleados en este libro,
con indicación de su función y del capítulo en que se desarrolla cada uno.

\begingroup
\small
\begin{longtable}{L{2.8cm} L{2.2cm} L{6cm} C{1.8cm}}
	% --- Título y encabezado de la primera página ---
	\caption{Diagramas -UML- empleados en el libro con su función y capítulo de desarrollo. Fuente: \textcite{Larman2004}.} \label{tab:diagramas-uml} \\
	\toprule
	\textbf{Diagrama} & \textbf{Familia} & \textbf{Función en el libro} & \textbf{Capítulo} \\
	\midrule
	\endfirsthead
	
	% --- Encabezado para las páginas siguientes ---
	\multicolumn{4}{c}{{\bfseries \tablename\ \thetable{} -- Continuación}} \\
	\toprule
	\textbf{Diagrama} & \textbf{Familia} & \textbf{Función en el libro} & \textbf{Capítulo} \\
	\midrule
	\endhead
	
	% --- Pie de tabla para páginas intermedias ---
	\midrule
	\multicolumn{4}{r}{{Continúa en la siguiente página...}} \\
	\bottomrule
	\endfoot
	
	% --- Pie de tabla final ---
	\bottomrule
	\endlastfoot
	
	% --- Datos de la tabla ---
	Clases (dominio)
	& Estructural
	& Modelar la jerarquía de entidades, atributos, relaciones y restricciones del dominio académico-administrativo de la \ud
	& 3 \\
	\midrule
	
	Casos de uso
	& Comportamiento
	& Especificar el comportamiento del sistema desde la perspectiva de los actores institucionales; derivar de las historias de usuario
	& 3 \\
	\midrule
	
	Secuencia
	& Comportamiento
	& Modelar la interacción dinámica entre objetos en los flujos de proceso más complejos (matrícula, evaluación con rúbrica)
	& 3 \\
	\midrule
	
	Estado
	& Comportamiento
	& Modelar el ciclo de vida de entidades con estados bien definidos (Estudiante, ProgramaAcademico)
	& 3 \\
	\midrule
	
	Paquetes
	& Estructural
	& Organizar las clases del dominio en los cinco módulos del sistema
	& 4, 5 \\
	\midrule
	
	Componentes
	& Estructural
	& Modelar la arquitectura del sistema y sus dependencias de despliegue
	& 6 \\
\end{longtable}
\endgroup

\plantuml\ se integra con el flujo de trabajo de compilación \LaTeX\ mediante
la herramienta \texttt{plantuml.jar}, que convierte los archivos \texttt{.puml}
en imágenes -SVG- o -PNG- que \LaTeX\ puede incluir con \verb|\includegraphics|.
En este libro, los diagramas se presentan como listados de código -PlantUML-
reproducibles: el lector puede copiar el código y generar la imagen de manera
local, lo que garantiza que los modelos sean verificables y versionables.
La convención adoptada para los identificadores -UML- del libro es:
\texttt{lst:plantuml-<tipo>-<subsistema>} para los listados
y \texttt{fig:uml-<tipo>-<subsistema>} para las figuras.

%% ─────────────────────────────────────────────
\section{Diagramas de casos de uso del sistema}
\label{sec:casos-uso}
%% ─────────────────────────────────────────────

Los diagramas de casos de uso se construyen a partir de las historias de
usuario del \textit{product backlog} (Sección~\ref{sec:historias-usuario})
y los requisitos funcionales (Sección~\ref{sec:req-funcionales}). Cada
actor institucional se asocia con los casos de uso que le conciernen,
organizados en tres subsistemas funcionales.

\subsection{Casos de uso del subsistema organizacional}
\label{subsec:uc-organizacional}

\begin{lstlisting}[style=plantumlStyle,
  caption={Diagrama de casos de uso del subsistema organizacional (PlantUML).},
  label={lst:plantuml-uc-org}]
@startuml uc_organizacional
skinparam actorStyle awesome
skinparam usecaseBackgroundColor #F0F4FF
skinparam usecaseBorderColor #00467F
skinparam arrowColor #00467F

left to right direction

actor "Rector" as rector
actor "Decano" as decano
actor "Director de Escuela" as dir_escuela
actor "Personal\nadministrativo" as admin

rectangle "Subsistema Organizacional" {
  usecase "Registrar Facultad" as RF_FAC
  usecase "Modificar datos de Facultad" as MOD_FAC
  usecase "Registrar Escuela" as RF_ESC
  usecase "Adscribir docente\na Escuela" as ADSC_DOC
  usecase "Registrar Centro" as RF_CEN
  usecase "Registrar Instituto" as RF_INST
  usecase "Registrar CABA" as RF_CABA
  usecase "Consultar estructura\norganizacional" as CONS_ORG
  usecase "Generar reporte de\ncomposicion institucional" as REP_COMP
  usecase "Gestionar directivo\nde unidad" as GEST_DIR
}

rector --> RF_FAC
rector --> MOD_FAC
rector --> CONS_ORG
rector --> REP_COMP

decano --> RF_ESC
decano --> RF_CEN
decano --> RF_INST
decano --> RF_CABA
decano --> CONS_ORG
decano --> GEST_DIR

dir_escuela --> ADSC_DOC
dir_escuela --> CONS_ORG

admin --> CONS_ORG
admin --> REP_COMP

RF_ESC ..> RF_FAC : <<include>>
RF_CABA ..> RF_ESC : <<include>>
ADSC_DOC ..> RF_ESC : <<include>>

@enduml
\end{lstlisting}

Los casos de uso del subsistema organizacional reflejan directamente las
restricciones del \acuerdo: registrar una Escuela presupone que ya existe
una Facultad a la que adscribirla (\textit{include}); adscribir un docente
presupone que la Escuela de destino ya existe. La relación de inclusión
expresa estas dependencias normativas de forma explícita en el modelo.

\subsection{Casos de uso del subsistema académico}
\label{subsec:uc-academico}

\begin{lstlisting}[style=plantumlStyle,
  caption={Diagrama de casos de uso del subsistema académico (PlantUML).},
  label={lst:plantuml-uc-acad}]
@startuml uc_academico
skinparam actorStyle awesome
skinparam usecaseBackgroundColor #F0F4FF
skinparam usecaseBorderColor #00467F
skinparam arrowColor #00467F

left to right direction

actor "Director de Programa" as dir_prog
actor "Docente" as docente
actor "Estudiante" as estudiante
actor "Personal\nadministrativo" as admin

rectangle "Subsistema Academico" {
  usecase "Registrar Programa\nAcademico" as RF_PROG
  usecase "Gestionar Plan\nde Estudios" as GEST_PLAN
  usecase "Registrar Espacio\nAcademico" as RF_ESP
  usecase "Asignar docente a\nespacio academico" as ASIG_DOC
  usecase "Matricular Estudiante" as MATR
  usecase "Consultar oferta\nacademica" as CONS_OFE
  usecase "Actualizar PEP" as ACT_PEP
  usecase "Generar reporte\nde cobertura" as REP_COB
  usecase "Consultar historial\nacademico" as CONS_HIST
}

dir_prog --> RF_PROG
dir_prog --> GEST_PLAN
dir_prog --> RF_ESP
dir_prog --> ACT_PEP
dir_prog --> REP_COB

docente --> ASIG_DOC
docente --> CONS_OFE

estudiante --> CONS_HIST
estudiante --> CONS_OFE

admin --> MATR
admin --> CONS_OFE

GEST_PLAN ..> RF_PROG : <<include>>
RF_ESP ..> GEST_PLAN : <<include>>
ASIG_DOC ..> RF_ESP : <<include>>
MATR ..> RF_PROG : <<include>>

ACT_PEP .> GEST_PLAN : <<extend>>

@enduml
\end{lstlisting}

\subsection{Casos de uso del subsistema de evaluación}
\label{subsec:uc-evaluacion}

\begin{lstlisting}[style=plantumlStyle,
  caption={Diagrama de casos de uso del subsistema de evaluación y -PFA-
    (PlantUML).},
  label={lst:plantuml-uc-eval}]
@startuml uc_evaluacion
skinparam actorStyle awesome
skinparam usecaseBackgroundColor #F0F4FF
skinparam usecaseBorderColor #00467F
skinparam arrowColor #00467F

left to right direction

actor "Director de Programa" as dir_prog
actor "Docente" as docente
actor "Estudiante" as estudiante
actor "Coordinador CABA" as coord

rectangle "Subsistema de Evaluacion y PFA" {
  usecase "Registrar RA\ndel programa" as RF_RA
  usecase "Registrar PFA de\nespacio academico" as RF_PFA
  usecase "Definir Rubrica\nde evaluacion" as DEF_RUB
  usecase "Registrar Evidencia\nde Aprendizaje" as RF_EVID
  usecase "Calcular avance del\nestudiante frente a RA" as CALC_AV
  usecase "Generar reporte de\nseguimiento PFA" as REP_PFA
  usecase "Consultar avance\npropio" as CONS_AV
}

dir_prog --> RF_RA
dir_prog --> RF_PFA
dir_prog --> REP_PFA

docente --> DEF_RUB
docente --> RF_EVID
docente --> CALC_AV

coord --> RF_PFA
coord --> REP_PFA

estudiante --> CONS_AV

RF_EVID ..> DEF_RUB : <<include>>
CALC_AV ..> RF_EVID : <<include>>
REP_PFA ..> RF_PFA : <<include>>
CONS_AV ..> CALC_AV : <<include>>

@enduml
\end{lstlisting}

\subsubsection*{Descripción expandida: caso de uso «Registrar Evidencia de Aprendizaje»}

Este caso de uso concentra la mayor complejidad del subsistema de evaluación,
toda vez que involucra a docente, rúbrica, evidencia y cálculo de avance en
una secuencia interdependiente.

\textbf{Actor principal:} Docente.\\
\textbf{Actores secundarios:} Estudiante (receptor de la calificación),
Director de Programa (consulta el avance).

\textbf{Precondición:} el espacio académico tiene al menos una rúbrica
definida (ver caso de uso «Definir Rúbrica»); el estudiante está matriculado
en el espacio académico.

\textbf{Flujo principal:}
\begin{enumerate}[label=\arabic*.]
  \item El docente accede al espacio académico asignado y selecciona el
    estudiante a evaluar.
  \item El sistema presenta las rúbricas disponibles para el espacio académico.
  \item El docente selecciona la rúbrica correspondiente a la actividad
    formativa.
  \item El sistema muestra los criterios de la rúbrica con los niveles
    de desempeño.
  \item El docente selecciona el nivel de desempeño para cada criterio
    y, opcionalmente, adjunta un archivo de evidencia.
  \item El sistema calcula la calificación según la ponderación de la rúbrica.
  \item El docente confirma el registro y el sistema almacena la evidencia
    con fecha, hora y actor responsable (RNF-003).
  \item El sistema actualiza el cálculo de avance del estudiante frente
    a los RA del programa.
\end{enumerate}

\textbf{Flujo alternativo A} (el docente detecta un error en la calificación):
el docente puede modificar el registro dentro de las 48 horas siguientes.
Pasado ese plazo, solo el Director de Programa puede modificarlo, dejando
trazabilidad del cambio.

\textbf{Postcondición:} la evidencia queda registrada, la calificación
almacenada y el avance del estudiante actualizado.

%% ─────────────────────────────────────────────
\section{Diagrama de clases del dominio}
\label{sec:clases-dominio}
%% ─────────────────────────────────────────────

El diagrama de clases del dominio se presenta en cuatro vistas complementarias.
La representación del dominio completo en un único diagrama resultaría
inmanejable ---y didácticamente contraproducente---, toda vez que la
complejidad de la jerarquía de unidades académicas, actores, programas
y evaluación supera con creces lo que un diagrama de una sola página puede
comunicar con claridad. Cada vista se acompaña de una justificación
de las decisiones de diseño más relevantes.

\subsection{Jerarquía de unidades académicas}
\label{subsec:jerarquia-unidades}

\begin{lstlisting}[style=plantumlStyle,
  caption={Jerarquía completa de unidades académicas de la \ud\ (PlantUML).},
  label={lst:plantuml-clases-org}]
@startuml clases_organizacional
skinparam classAttributeIconSize 0
skinparam classFontStyle bold
skinparam classBackgroundColor #F8F8FF
skinparam classBorderColor #00467F
skinparam ArrowColor #00467F

abstract class UnidadAcademica {
  # nombre : String
  # codigo : String
  # fechaCreacion : LocalDate
  # estado : EstadoUnidad
  + getNombre() : String
  + getCodigo() : String
  + getEstado() : EstadoUnidad
  + {abstract} getTipo() : String
  + {abstract} getDirector() : String
}

class Universidad {
  - nombre : String
  - nit : String
  - mision : String
  + getFacultades() : List<Facultad>
  + agregarFacultad(f : Facultad) : void
}

class Facultad {
  - decano : String
  - codigoFacultad : String
  + getDecano() : String
  + getEscuelas() : List<Escuela>
  + agregarEscuela(e : Escuela) : void
  + getCentros() : List<Centro>
  + getInstitutos() : List<Instituto>
  + getTipo() : String
  + getDirector() : String
}

class Escuela {
  - directorEscuela : String
  + getDirectorEscuela() : String
  + getDocentesAdscritos() : List<Docente>
  + adscribirDocente(d : Docente) : void
  + getCABAs() : List<CABA>
  + getTipo() : String
  + getDirector() : String
}

class Centro {
  - director : String
  - mision : String
  + getTipo() : String
  + getDirector() : String
}

class Instituto {
  - director : String
  - areaTematica : String
  + getTipo() : String
  + getDirector() : String
}

class CABA {
  - coordinador : String
  + getCoordinador() : String
  + getAreaDeFormacion() : AreaDeFormacion
  + getTipo() : String
  + getDirector() : String
}

class AreaDeFormacion {
  - nombre : String
  - codigo : String
  + getProgramas() : List<ProgramaAcademico>
  + agregarPrograma(p : ProgramaAcademico) : void
}

class ProgramaAcademico {
  - nombre : String
  - codigoSNIES : String
  - nivel : NivelAcademico
  - estado : EstadoPrograma
  - fechaRegistroCalificado : LocalDate
  + getNombre() : String
  + getNivel() : NivelAcademico
  + getPlanDeEstudios() : PlanDeEstudios
}

class PlanDeEstudios {
  - version : String
  - totalCreditos : int
  - fechaAprobacion : LocalDate
  + getEspacios() : List<EspacioAcademico>
  + agregarEspacio(e : EspacioAcademico) : void
  + getTotalCreditos() : int
}

class EspacioAcademico {
  - nombre : String
  - codigo : String
  - creditos : int
  - nivelFormacion : String
  + getNombre() : String
  + getCreditos() : int
  + getDocenteAsignado() : Docente
}

enum EstadoUnidad {
  ACTIVA
  INACTIVA
  EN_CREACION
}

enum EstadoPrograma {
  EN_DISENO
  EN_TRAMITE
  ACTIVO
  EN_RENOVACION
  CERRADO
}

enum NivelAcademico {
  PREGRADO
  ESPECIALIZACION
  MAESTRIA
  DOCTORADO
}

UnidadAcademica <|-- Facultad
UnidadAcademica <|-- Escuela
UnidadAcademica <|-- Centro
UnidadAcademica <|-- Instituto

Universidad "1" *-- "1..*" Facultad : contiene
Facultad "1" *-- "0..*" Escuela : agrupa
Facultad "1" o-- "0..*" Centro : contiene
Facultad "1" o-- "0..*" Instituto : contiene
Escuela "1" o-- "0..*" CABA : coordina
Facultad "1" o-- "1..*" AreaDeFormacion : organiza
AreaDeFormacion "1" --> "1..*" ProgramaAcademico : ofrece
ProgramaAcademico "1" *-- "1" PlanDeEstudios : define
PlanDeEstudios "1" *-- "1..*" EspacioAcademico : contiene

UnidadAcademica --> EstadoUnidad
ProgramaAcademico --> EstadoPrograma
ProgramaAcademico --> NivelAcademico

@enduml
\end{lstlisting}

\textbf{Justificación de las decisiones de diseño.} La clase
\texttt{UnidadAcademica} se modela como abstracta porque ningún objeto del
sistema es ``una unidad académica genérica'': siempre se trata de una Facultad,
una Escuela, un Centro o un Instituto concreto. El método \texttt{getTipo()}
abstracto fuerza a cada subclase a declarar su tipo, lo que permite escribir
código polimórfico que genere reportes sin conocer el tipo concreto de cada
unidad. La relación entre \texttt{Facultad} y \texttt{Escuela} es de
composición ---representada con el rombo relleno--- porque el \acuerdo\
establece que las Escuelas pertenecen a una Facultad y no tienen existencia
autónoma fuera de ella. La relación entre \texttt{Escuela} y \texttt{CABA}
es de agregación, puesto que el estatuto sugiere que los CABAS son mecanismos
de coordinación que articulan docencia e investigación, no subunidades
con existencia completamente dependiente de la Escuela.\footnote{%
  La distinción entre composición y agregación ---rombo relleno vs.\ rombo
  vacío en la notación -UML--- no siempre es inmediata a partir del texto
  del estatuto. En la experiencia docente acumulada entre 2020 y 2025 en
  la Facultad de Ingeniería de la \ud, sede Macarena, este es uno de los
  puntos donde los estudiantes más frecuentemente consultan al asistente
  computacional y aceptan su respuesta sin verificar si la semántica es
  coherente con el dominio. La regla práctica es: composición cuando la
  destrucción del contenedor implica la destrucción del contenido;\\aggregación
  cuando el contenido puede existir de forma independiente.}

\subsection{Jerarquía de actores institucionales}
\label{subsec:jerarquia-actores}

\begin{lstlisting}[style=plantumlStyle,
  caption={Jerarquía de actores institucionales (PlantUML).
    Clase abstracta \texttt{Persona} y sus subclases según el Art.\ 8
    del \acuerdo.},
  label={lst:plantuml-clases-actores}]
@startuml clases_actores
skinparam classAttributeIconSize 0
skinparam classFontStyle bold
skinparam classBackgroundColor #F8F8FF
skinparam classBorderColor #00467F
skinparam ArrowColor #00467F

abstract class Persona {
  # nombres : String
  # apellidos : String
  # identificacion : String
  # tipoId : TipoIdentificacion
  # correoInstitucional : String
  + getNombreCompleto() : String
  + getIdentificacion() : String
  + {abstract} getRol() : String
}

class Docente {
  - tipoVinculacion : TipoDocente
  - fechaVinculacion : LocalDate
  - escuelaAdscrita : Escuela
  - especialidades : List<String>
  + getTipoVinculacion() : TipoDocente
  + getEscuelaAdscrita() : Escuela
  + getEspaciosAsignados() : List<EspacioAcademico>
  + getRol() : String
}

class Estudiante {
  - codigoEstudiante : String
  - estado : EstadoEstudiante
  - programaMatriculado : ProgramaAcademico
  - creditosAcumulados : int
  + getCodigoEstudiante() : String
  + getEstado() : EstadoEstudiante
  + getProgramaMatriculado() : ProgramaAcademico
  + getAvanceRA() : double
  + getRol() : String
}

class PersonalAdministrativo {
  - cargo : String
  - dependencia : String
  - nivelCargo : NivelCargo
  + getCargo() : String
  + getDependencia() : String
  + getRol() : String
}

class Egresado {
  - codigoEstudianteOriginal : String
  - programaGraduado : ProgramaAcademico
  - fechaGrado : LocalDate
  - tituloObtenido : String
  + getProgramaGraduado() : ProgramaAcademico
  + getTituloObtenido() : String
  + getRol() : String
}

class Matricula {
  - periodo : String
  - fechaMatricula : LocalDate
  - estado : EstadoMatricula
  - creditosInscritos : int
  + getPeriodo() : String
  + getEstado() : EstadoMatricula
  + cancelar() : void
}

enum TipoDocente {
  PLANTA
  OCASIONAL
  HORA_CATEDRA
  VISITANTE
  EXPERTO
}

enum EstadoEstudiante {
  ASPIRANTE
  ADMITIDO
  ACTIVO
  SUSPENSION
  EN_GRADO
  GRADUADO
  RETIRADO
}

enum EstadoMatricula {
  ACTIVA
  CANCELADA
  BLOQUEADA
  EN_GRADO
}

enum TipoIdentificacion {
  CC
  TI
  CE
  PASAPORTE
}

enum NivelCargo {
  DIRECTIVO
  ASESOR
  PROFESIONAL
  TECNICO
  ASISTENCIAL
}

Persona <|-- Docente
Persona <|-- Estudiante
Persona <|-- PersonalAdministrativo
Persona <|-- Egresado

Docente --> TipoDocente : tipoVinculacion
Docente --> Escuela : adscrito a
Estudiante --> EstadoEstudiante
Matricula --> EstadoMatricula
Persona --> TipoIdentificacion

Estudiante "1" --> "1..*" Matricula : tiene
Matricula --> ProgramaAcademico : en

@enduml
\end{lstlisting}

\textbf{Justificación.} La clase abstracta \texttt{Persona} agrupa los
atributos de identificación comunes a todos los miembros de la comunidad
universitaria: nombre, identificación y correo institucional. El método
\texttt{getRol()} abstracto fuerza a cada subclase a declarar su rol en
el sistema, lo que permite construir componentes de -UI- genéricos que
muestren el rol del usuario autenticado sin un bloque de condicionales.

La decisión de no modelar subclases de \texttt{Docente} para los cinco tipos
de vinculación ---y utilizar en cambio el enumerado \texttt{TipoDocente}---
responde a que, en el dominio de la \ud\ según el \acuerdo, los cinco tipos
no difieren en comportamiento sino en atributos: el tipo de contrato,
la fecha de vinculación y algunos derechos administrativos. Si en una fase
posterior del sistema se requiriese modelar la liquidación de honorarios
---proceso cuyas reglas difieren sustancialmente por tipo de vinculación---,
la conversión del enumerado en una jerarquía de subclases podría justificarse.
Por ahora, la solución más simple es la más adecuada.\footnote{%
  El principio \textsc{YAGNI} (\textit{You Aren't Gonna Need It}), aunque
  proveniente de las metodologías ágiles y no del -UML- clásico, aplica
  directamente a esta decisión: modelar las subclases de \texttt{Docente}
  ahora, antes de que exista un requisito funcional que las justifique,
  introduce complejidad sin beneficio observable. \textcite{Fowler2018}
  cataloga este antipatrón como \textit{Speculative Generality}: la
  generalización prematura que ninguna funcionalidad real solicita.}

\subsection{Programas, planes de estudio y espacios académicos}
\label{subsec:modelo-programas}

\begin{lstlisting}[style=plantumlStyle,
  caption={Modelo de programas académicos, planes de estudio y espacios
    académicos (PlantUML).},
  label={lst:plantuml-clases-acad}]
@startuml clases_academico
skinparam classAttributeIconSize 0
skinparam classFontStyle bold
skinparam classBackgroundColor #F8F8FF
skinparam classBorderColor #00467F
skinparam ArrowColor #00467F

class ProgramaAcademico {
  - nombre : String
  - codigoSNIES : String
  - nivel : NivelAcademico
  - estado : EstadoPrograma
  - fechaRegistro : LocalDate
  + getNombre() : String
  + getPlanVigente() : PlanDeEstudios
  + getResultadosAprendizaje() : List<ResultadoDeAprendizaje>
}

class PlanDeEstudios {
  - version : String
  - totalCreditos : int
  - fechaAprobacion : LocalDate
  - organoAprobador : String
  + getVersion() : String
  + getTotalCreditos() : int
  + getEspaciosPorSemestre(sem : int) : List<EspacioAcademico>
}

class EspacioAcademico {
  - nombre : String
  - codigo : String
  - creditos : int
  - semestre : int
  - nivelFormacion : String
  + getCreditos() : int
  + getPrerrequisitos() : List<EspacioAcademico>
  + getCorequisitos() : List<EspacioAcademico>
  + getPFA() : List<PropositoDeFormacion>
}

class ProyectoEducativoDePrograma {
  - version : String
  - fechaAprobacion : LocalDate
  - perfilEgresado : String
  - perfilIngreso : String
  + getResultadosAprendizaje() : List<ResultadoDeAprendizaje>
  + getPlanVigente() : PlanDeEstudios
}

class ResultadoDeAprendizaje {
  - codigo : String
  - descripcion : String
  - verboBloom : String
  - nivel : NivelBloom
  + getCodigo() : String
  + getDescripcion() : String
  + getNivel() : NivelBloom
}

class PropositoDeFormacion {
  - ambito : AmbitoPFA
  - descripcion : String
  + getAmbito() : AmbitoPFA
  + getResultadosVinculados() : List<ResultadoDeAprendizaje>
}

enum AmbitoPFA {
  ONTOLOGICO
  EPISTEMOLOGICO
  CONTEXTUAL
}

enum NivelBloom {
  RECORDAR
  COMPRENDER
  APLICAR
  ANALIZAR
  EVALUAR
  CREAR
}

ProgramaAcademico "1" *-- "1..*" PlanDeEstudios : versiona
ProgramaAcademico "1" *-- "1" ProyectoEducativoDePrograma : define
PlanDeEstudios "1" *-- "1..*" EspacioAcademico : contiene
EspacioAcademico "0..*" o-- "0..*" EspacioAcademico : prerrequisito de >
EspacioAcademico "1" *-- "0..*" PropositoDeFormacion : define
PropositoDeFormacion "1..*" --> "1..*" ResultadoDeAprendizaje : vincula
ProgramaAcademico "1" --> "1..*" ResultadoDeAprendizaje : declara
PropositoDeFormacion --> AmbitoPFA
ResultadoDeAprendizaje --> NivelBloom

@enduml
\end{lstlisting}

\textbf{Justificación.} La relación reflexiva de \texttt{EspacioAcademico}
con sí mismo modela la restricción de prerrequisitos: un espacio académico
puede tener cero o más prerrequisitos, que son otros espacios académicos
del mismo plan de estudios. La multiplicidad 0..* en ambos extremos expresa
que algunos espacios no tienen prerrequisitos (electivas de primer semestre)
y que un espacio puede ser prerrequisito de varios otros.

La existencia de múltiples versiones del \texttt{PlanDeEstudios} para un
mismo \texttt{ProgramaAcademico} ---relación de composición uno a muchos---
refleja la realidad normativa: los programas reforman sus planes de estudios
periódicamente, pero los estudiantes que iniciaron bajo un plan anterior
conservan el derecho a terminarlo bajo las condiciones originales. El sistema
debe mantener el histórico de versiones y asociar cada estudiante a la versión
del plan bajo la cual se matriculó.

\subsection{Modelo de evaluación, -PFA- y rúbricas}
\label{subsec:modelo-evaluacion}

\begin{lstlisting}[style=plantumlStyle,
  caption={Modelo de evaluación, rúbricas y evidencias (PlantUML).},
  label={lst:plantuml-clases-eval}]
@startuml clases_evaluacion
skinparam classAttributeIconSize 0
skinparam classFontStyle bold
skinparam classBackgroundColor #F8F8FF
skinparam classBorderColor #00467F
skinparam ArrowColor #00467F

class Rubrica {
  - nombre : String
  - descripcion : String
  - version : String
  - ponderacionTotal : double
  + getNombre() : String
  + getCriterios() : List<Criterio>
  + calcularCalificacion(niveles : Map<Criterio, NivelDesempeno>) : double
}

class Criterio {
  - nombre : String
  - descripcion : String
  - ponderacion : double
  + getNombre() : String
  + getPonderacion() : double
  + getNivelesDesempeno() : List<NivelDesempeno>
}

class NivelDesempeno {
  - nombre : String
  - descripcion : String
  - valorNumerico : double
  + getNombre() : String
  + getValorNumerico() : double
}

class Evidencia {
  - titulo : String
  - descripcion : String
  - tipo : TipoEvidencia
  - fechaRegistro : LocalDateTime
  - archivoRuta : String
  - docente : Docente
  - estudiante : Estudiante
  + getTitulo() : String
  + getFechaRegistro() : LocalDateTime
  + getCriteriosEvaluados() : List<CriterioEvaluado>
}

class CriterioEvaluado {
  - criterio : Criterio
  - nivelObtenido : NivelDesempeno
  - comentario : String
  + getCalificacionPonderada() : double
}

class Calificacion {
  - valor : double
  - escala : String
  - fecha : LocalDate
  - periodo : String
  + getValor() : double
  + esAprobatorio() : boolean
}

class ActividadFormativa {
  - nombre : String
  - descripcion : String
  - tipo : TipoActividad
  - porcentaje : double
  + getNombre() : String
  + getPorcentaje() : double
}

enum TipoEvidencia {
  TRABAJO_ESCRITO
  PRESENTACION
  PROYECTO
  EXAMEN
  PARTICIPACION
  OTRO
}

enum TipoActividad {
  TALLER
  PROYECTO
  EXAMEN
  EXPOSICION
  PRACTICA
  INVESTIGACION
}

Rubrica "1" *-- "1..*" Criterio : contiene
Criterio "1" *-- "2..4" NivelDesempeno : define
Evidencia "1" *-- "1..*" CriterioEvaluado : registra
CriterioEvaluado --> Criterio
CriterioEvaluado --> NivelDesempeno
Evidencia --> Docente : registrada por
Evidencia --> Estudiante : evaluado
EspacioAcademico "1" --> "1..*" Rubrica : emplea
EspacioAcademico "1" --> "0..*" ActividadFormativa : incluye
Evidencia "1" --> "1" Calificacion : genera
Evidencia --> TipoEvidencia
ActividadFormativa --> TipoActividad

@enduml
\end{lstlisting}

\textbf{Justificación.} La separación entre \texttt{Rubrica},
\texttt{Criterio} y \texttt{NivelDesempeno} modela la estructura real
de las rúbricas de evaluación: una rúbrica tiene varios criterios, y
cada criterio tiene entre dos y cuatro niveles de desempeño que varían
en descripción y valor numérico. La clase \texttt{CriterioEvaluado} es
una clase asociativa que registra, para cada evidencia concreta, qué nivel
de desempeño obtuvo el estudiante en cada criterio. Este diseño permite
recalcular la calificación en cualquier momento a partir de los registros
originales y mantiene la trazabilidad completa del proceso evaluativo.

%% ─────────────────────────────────────────────
\section{Diagramas de secuencia representativos}
\label{sec:secuencia}
%% ─────────────────────────────────────────────

Los diagramas de secuencia modelan la interacción dinámica entre objetos
del dominio en los flujos de mayor relevancia. Se presentan dos: la
matrícula de un estudiante y la evaluación con rúbrica. Cada diagrama
se acompaña de una descripción de las responsabilidades asignadas a
cada participante.

\subsection*{Secuencia: Matricular Estudiante}

\begin{lstlisting}[style=plantumlStyle,
  caption={Diagrama de secuencia: matricular estudiante (PlantUML).},
  label={lst:plantuml-seq-matricula}]
@startuml seq_matricula
skinparam sequenceArrowColor #00467F
skinparam participantBackgroundColor #F0F4FF
skinparam participantBorderColor #00467F
skinparam lifelineBorderColor #00467F

actor "Personal\nAdministrativo" as admin
participant "SistemaMatricula\n:ControladorMatricula" as ctrl
participant ":Estudiante" as est
participant ":ProgramaAcademico" as prog
participant ":Matricula" as mat
participant ":RepositorioMatricula" as repo

admin -> ctrl : registrarMatricula(datosEstudiante, codigoProg, periodo)
activate ctrl

ctrl -> prog : verificarCuposDisponibles(periodo)
activate prog
prog --> ctrl : cuposDisponibles : boolean
deactivate prog

alt cuposDisponibles == true
  ctrl -> est : create(datosEstudiante)
  activate est
  est --> ctrl : estudiante : Estudiante
  deactivate est

  ctrl -> mat : create(estudiante, programa, periodo)
  activate mat
  mat --> ctrl : matricula : Matricula
  deactivate mat

  ctrl -> repo : guardar(matricula)
  activate repo
  repo --> ctrl : confirmacion : boolean
  deactivate repo

  ctrl -> prog : reducirCupo(periodo)
  activate prog
  prog --> ctrl : ok
  deactivate prog

  ctrl --> admin : matriculaRegistrada(matricula.getCodigo())
else cuposDisponibles == false
  ctrl --> admin : error("Sin cupos disponibles en el periodo")
end

deactivate ctrl
@enduml
\end{lstlisting}

El controlador \texttt{ControladorMatricula} concentra la lógica de
orquestación: verifica cupos, crea los objetos de dominio, persiste la
matrícula y actualiza el estado del programa. La separación entre el
controlador y las clases de dominio respeta el principio de responsabilidad
única (\textit{Single Responsibility Principle}): las clases de dominio
modelan conceptos del negocio, no lógica de transacción.

\subsection*{Secuencia: Evaluar con Rúbrica}

\begin{lstlisting}[style=plantumlStyle,
  caption={Diagrama de secuencia: registrar evidencia y evaluar con
    rúbrica (PlantUML).},
  label={lst:plantuml-seq-rubrica}]
@startuml seq_rubrica
skinparam sequenceArrowColor #00467F
skinparam participantBackgroundColor #F0F4FF
skinparam participantBorderColor #00467F

actor "Docente" as docente
participant "ControladorEvaluacion" as ctrl
participant ":EspacioAcademico" as espacio
participant ":Rubrica" as rubrica
participant ":Evidencia" as evid
participant ":Calificacion" as cal
participant ":RepositorioEvaluacion" as repo

docente -> ctrl : iniciarEvaluacion(idEspacio, idEstudiante)
activate ctrl

ctrl -> espacio : getRubricas()
activate espacio
espacio --> ctrl : List<Rubrica>
deactivate espacio

ctrl --> docente : presentarRubricas(List<Rubrica>)

docente -> ctrl : seleccionarRubrica(idRubrica)
ctrl -> rubrica : getCriterios()
activate rubrica
rubrica --> ctrl : List<Criterio>
deactivate rubrica

ctrl --> docente : presentarCriterios(List<Criterio>)

docente -> ctrl : registrarNiveles(Map<Criterio, NivelDesempeno>, archivo)
ctrl -> evid : create(docente, estudiante, rubrica, nivelesObtenidos, archivo)
activate evid
evid --> ctrl : evidencia : Evidencia
deactivate evid

ctrl -> rubrica : calcularCalificacion(nivelesObtenidos)
activate rubrica
rubrica --> ctrl : valorCalificacion : double
deactivate rubrica

ctrl -> cal : create(valorCalificacion, periodo)
activate cal
cal --> ctrl : calificacion : Calificacion
deactivate cal

ctrl -> repo : guardar(evidencia, calificacion)
activate repo
repo --> ctrl : confirmacion
deactivate repo

ctrl --> docente : evaluacionRegistrada(calificacion.getValor())
deactivate ctrl
@enduml
\end{lstlisting}

La secuencia de evaluación con rúbrica evidencia la ventaja del diseño
encapsulado: el docente no necesita conocer el algoritmo de cálculo de
la calificación ---ese conocimiento está en \texttt{Rubrica.calcularCalificacion()}---
ni los detalles de persistencia ---responsabilidad del repositorio---. Solo
provee los datos que le corresponden como actor del proceso.

%% ─────────────────────────────────────────────
\section{Diagramas de estado para procesos clave}
\label{sec:estados}
%% ─────────────────────────────────────────────

Los diagramas de estado modelan el ciclo de vida de las entidades del
dominio con mayor variación de estado. Dos resultan determinantes para
el sistema de información de la \ud: el ciclo de vida del
\texttt{Estudiante} y el del \texttt{ProgramaAcademico}.

\subsection*{Estado: Ciclo de vida del Estudiante}

\begin{lstlisting}[style=plantumlStyle,
  caption={Diagrama de estado: ciclo de vida del \texttt{Estudiante}
    (PlantUML).},
  label={lst:plantuml-estado-estudiante}]
@startuml estado_estudiante
skinparam stateBackgroundColor #F0F4FF
skinparam stateBorderColor #00467F
skinparam arrowColor #00467F

[*] --> Aspirante : Solicitud de admisión

Aspirante --> Admitido : Aprobación del proceso de admisión
Aspirante --> [*] : Rechazo o desistimiento

Admitido --> Activo : Pago de matrícula y\nregistro académico
Admitido --> [*] : No pago en plazo establecido

Activo --> Suspension : Sanción disciplinaria\no abandono por período
Activo --> EnProcesodeGrado : Aprobación de todos\nlos créditos requeridos
Activo --> Retirado : Retiro voluntario\no definitivo
Activo --> Activo : Renovación de matrícula\npor período

Suspension --> Activo : Levantamiento de sanción\no reingreso aprobado
Suspension --> Retirado : Vencimiento del plazo\nde reingreso

EnProcesodeGrado --> Graduado : Grado conferido en\nceremonia oficial
EnProcesodeGrado --> Activo : Pendiente de requisitos\nde grado

Retirado --> [*] : Estado terminal
Graduado --> [*] : Estado terminal\n(pasa a Egresado)

state "En Proceso de Grado" as EnProcesodeGrado
@enduml
\end{lstlisting}

Los estados \texttt{Graduado} y \texttt{Retirado} son terminales: una vez
que un estudiante alcanza cualquiera de ellos, no puede regresar al estado
\texttt{Activo}. Esta restricción debe reflejarse en el código mediante
una excepción de dominio ---por ejemplo, \texttt{TransicionEstadoInvalidaException}---
que impida la transición ilegal. El patrón \textit{State} \cite{Gamma1994}
es el candidato natural para implementar este comportamiento en \java.\footnote{%
  En \rust\, el patrón \textit{State} se puede implementar mediante enumerados
  con variantes que llevan datos (\textit{enums with data}), lo que resulta
  más expresivo que la solución basada en interfaces de \java. La comparación
  detallada se desarrolla en el Capítulo correspondiente a la implementación
  en \rust.}

\subsection*{Estado: Ciclo de vida del Programa Académico}

\begin{lstlisting}[style=plantumlStyle,
  caption={Diagrama de estado: ciclo de vida del \texttt{ProgramaAcademico}
    (PlantUML).},
  label={lst:plantuml-estado-programa}]
@startuml estado_programa
skinparam stateBackgroundColor #F0F4FF
skinparam stateBorderColor #00467F
skinparam arrowColor #00467F

[*] --> EnDiseno : Decisión del Consejo de Facultad\nde crear el programa

EnDiseno --> EnTramite : Radicación del documento\nante el Consejo Académico
EnDiseno --> EnDiseno : Revisión y ajuste\ndel PEP

EnTramite --> RegistroCalificado : Resolución del MEN\notorgando registro
EnTramite --> EnDiseno : Concepto desfavorable;\nse requieren ajustes

RegistroCalificado --> Activo : Apertura de primera\ncohorte con matriculados
RegistroCalificado --> Cerrado : Decisión institucional\nde no abrir cohortes

Activo --> EnRenovacion : Inicio del proceso de\nrenovación del registro calificado
Activo --> Cerrado : Decisión del Consejo Superior\nde cerrar el programa

EnRenovacion --> Activo : Nueva resolución MEN\notorgando renovación
EnRenovacion --> Cerrado : No renovación;\nse inicia proceso de cierre

Cerrado --> [*] : Estado terminal

state "Registro Calificado" as RegistroCalificado
state "En Renovacion" as EnRenovacion
state "En Tramite" as EnTramite
@enduml
\end{lstlisting}

El ciclo de vida del \texttt{ProgramaAcademico} refleja la normativa del
Decreto~1330 de 2019 \autocite{Decreto1330_2019}: un programa no puede
tener estudiantes matriculados sin registro calificado vigente. El sistema
debe validar esta restricción al intentar matricular un estudiante en un
programa en estado \texttt{EnDiseno}, \texttt{EnTramite} o \texttt{Cerrado}.
La restricción es un invariante del dominio que el sistema de información
debe hacer cumplir.

%% ─────────────────────────────────────────────
\section{Comparativa \java\ vs.\ \rust: expresividad en modelado de dominio}
\label{sec:java-rust-uml}
%% ─────────────────────────────────────────────

La traducción de los diagramas de clases del dominio a código difiere
de manera notable entre \java\ y \rust. Tres aspectos ilustran el contraste.

\textbf{Clases abstractas.} La clase \texttt{UnidadAcademica} se mapea
directamente a \texttt{abstract class} en \java. En \rust\ no existe
el concepto de clase abstracta; la solución sintática es definir un
\textit{trait} \texttt{UnidadAcademica} con los métodos comunes e implementarlo
en cada \textit{struct} (\texttt{Facultad}, \texttt{Escuela}, etc.)
\autocite{Klabnik2023}. El \textit{trait} resulta más flexible ---permite
implementación múltiple, ausente en \java--- pero menos expresivo para
el concepto de herencia de datos: los atributos comunes deben duplicarse
en cada \textit{struct} o gestionarse mediante composición.

\textbf{Enumerados con datos.} El enumerado \texttt{EstadoEstudiante}
se modela en \java\ como un \texttt{enum} simple con constantes. En \rust\,
los enumerados pueden llevar datos en sus variantes, lo que permite
modelar estados con transiciones condicionadas por datos concretos ---por
ejemplo, la variante \texttt{EnProcesodeGrado\{fechaInicio: NaiveDate\}}
que registra la fecha en que comenzó el proceso de grado. Esta capacidad
hace que el patrón \textit{State} sea más natural en \rust\ que en \java.

\textbf{Relaciones de composición y propiedad.} La relación de composición
entre \texttt{PlanDeEstudios} y \texttt{EspacioAcademico} ---en la que el
plan es propietario del espacio y el espacio no tiene existencia independiente---
se expresa en \rust\ de forma directa mediante el sistema de \textit{ownership}:
\texttt{struct PlanDeEstudios \{ espacios: Vec<EspacioAcademico> \}} indica
que el \textit{struct} posee los espacios. En \java\, la misma semántica
debe documentarse mediante comentarios o convenciones de uso, puesto que
el lenguaje no distingue a nivel de tipos entre composición y agregación.

%% ─────────────────────────────────────────────
%% SPRINT 4
%% ─────────────────────────────────────────────
\section*{Sprint 4: Modelo de dominio inicial}
\label{sprint:4}
\addcontentsline{toc}{section}{Sprint 4: Modelo de dominio inicial}

\begin{tcolorbox}[sprintBox, title={Sprint 4 — Modelo de dominio inicial}]
  \textbf{Duración estimada:} 1 semana\\
  \textbf{Objetivo:} Construir el diagrama de clases del dominio para los
  subsistemas organizacional y de actores institucionales, empleando
  \plantuml\ y sustentando cada decisión de diseño en los requisitos
  derivados del \acuerdo.\\
  \textbf{Entregables:} (1) Diagrama de clases subsistema organizacional;
  (2) Diagrama de clases subsistema actores; (3) Documento de justificación
  de decisiones de diseño.
\end{tcolorbox}

\bigskip

\subsection*{Tabla del Sprint 4}

\begin{longtable}{L{3.5cm} L{4.5cm} L{4.5cm}}
  \sprintcabecera
  \endhead
  \bottomrule
  \caption{Entregables, características del programa y criterios de la
    rúbrica del Sprint 4.}
  \label{tab:sprint4}
  \endlastfoot
  Diagrama de clases del subsistema organizacional en \plantuml\
  (mínimo 8 clases con atributos, métodos y enumerados; relaciones
  de herencia, composición y asociación etiquetadas)
  & Aplicar los principios de abstracción, encapsulamiento y herencia
    al modelado de la jerarquía de unidades académicas; usar \plantuml\
    como herramienta de modelado versionable
  & El diagrama es coherente con el \acuerdo; las decisiones de herencia,
    composición y asociación se justifican mediante argumentación
    orientada a objetos con referencia a los artículos del estatuto;
    la sintaxis \plantuml\ es correcta y el diagrama compila \\
  \midrule
  Diagrama de clases del subsistema de actores en \plantuml\
  (mínimo 5 clases, incluyendo clase abstracta y enumerados de tipo
  de vinculación y estado)
  & Modelar la diversidad de actores institucionales mediante herencia
    y polimorfismo; justificar la decisión de enumerado vs.\ jerarquía
    de subclases para los tipos de docente
  & El diagrama distingue correctamente herencia de composición; los
    atributos de cada subclase son coherentes con el estatuto; la
    justificación de la decisión de diseño para \texttt{TipoDocente}
    está argumentada \\
  \midrule
  Documento de justificación de decisiones de diseño (mínimo 3 páginas)
  & Comunicar por escrito las decisiones de modelado, con argumentación
    orientada a objetos y referencia normativa; desarrollar la práctica
    de documentación de diseño
  & El documento justifica al menos 5 decisiones de diseño concretas
    con argumentación técnica y normativa; el vocabulario es coherente
    con los conceptos de -POO- desarrollados en el capítulo~1 \\
  \midrule
  Comparación con el modelo generado por asistente computacional
  (análisis crítico según el marco de 4 categorías)
  & Desarrollar pensamiento crítico sistemático frente al diseño
    orientado a objetos generado automáticamente
  & El análisis identifica al menos 3 diferencias conceptuales
    significativas; se aplican las categorías C1--C4 del marco de
    análisis; la argumentación está referenciada al dominio
    institucional \\
\end{longtable}

\subsection*{Actividad de reflexión: software generativo y modelado -UML-}
\label{act:ia-sprint4}

\begin{tcolorbox}[actividadIA,
  title={Actividad de reflexión (Sprint~4) —
         Software generativo y modelado -UML-}]

\textbf{Instrucción general.}
El estudiante elabora el diagrama de clases del subsistema organizacional
de manera individual, sin asistencia computacional. A continuación, solicita
al asistente generativo que genere un diagrama equivalente a partir de una
descripción en lenguaje natural del \acuerdo\ (proporcionada como contexto).
Se comparan las dos propuestas mediante el marco de análisis crítico.

\bigskip
\textbf{Preguntas orientadoras:}
\begin{enumerate}[label=\arabic*.]
  \item ¿El asistente genera código \plantuml\ sintácticamente correcto
        que compila sin errores? Si hay errores, identifíquelos y explique
        su probable causa.
  \item ¿El asistente distingue correctamente entre herencia y composición
        en el dominio institucional? Proporcione al menos un ejemplo concreto
        de error o imprecisión conceptual.
  \item ¿El asistente captura la restricción del \acuerdo\ según la cual
        todos los docentes están adscritos a una Escuela ---no a una
        Facultad---? Si no la captura, ¿cómo reformularía el \textit{prompt}
        para obtener un modelo más preciso?
  \item ¿El diagrama del asistente incluye enumerados para los estados
        y tipos? ¿Son coherentes con los valores definidos en el estatuto?
\end{enumerate}

\bigskip
\textbf{Rúbrica breve:}
\begin{itemize}
  \item \textbf{Producto} (60\%): corrección técnica del diagrama propio,
        coherencia con el \acuerdo, justificación de decisiones de diseño,
        compilación exitosa del código \plantuml.
  \item \textbf{Proceso} (40\%): calidad del análisis comparativo,
        identificación de errores conceptuales del asistente con
        argumentación normativa específica, reflexión sobre la mejora
        iterativa del \textit{prompt}.
\end{itemize}

\end{tcolorbox}

%% ─────────────────────────────────────────────
%% SPRINT 5
%% ─────────────────────────────────────────────
\section*{Sprint 5: Casos de uso ampliados y diagramas de secuencia}
\label{sprint:5}
\addcontentsline{toc}{section}{Sprint 5: Casos de uso ampliados y diagramas
  de secuencia}

\begin{tcolorbox}[sprintBox, title={Sprint 5 — Casos de uso ampliados y
  diagramas de secuencia}]
  \textbf{Duración estimada:} 1 semana\\
  \textbf{Objetivo:} Ampliar el diagrama de casos de uso a los tres
  subsistemas, especificar los flujos principal y alternativo de los
  casos de uso prioritarios, construir los diagramas de secuencia para
  los flujos principales y modelar el ciclo de vida de las entidades
  con estados.\\
  \textbf{Entregables:} (1) Diagramas CU completos (3 subsistemas);
  (2) Especificaciones textuales de 3 CU; (3) Diagramas de secuencia
  (matrícula + evaluación); (4) Diagramas de estado (Estudiante +
  ProgramaAcademico).
\end{tcolorbox}

\bigskip

\subsection*{Tabla del Sprint 5}

\begin{longtable}{L{3.5cm} L{4.5cm} L{4.5cm}}
  \sprintcabecera
  \endhead
  \bottomrule
  \caption{Entregables, características del programa y criterios de la
    rúbrica del Sprint 5.}
  \label{tab:sprint5}
  \endlastfoot
  Diagramas de casos de uso completos en \plantuml\ para los tres
  subsistemas (organizacional, académico, evaluación), con relaciones
  \textit{include} y \textit{extend} correctas
  & Modelar el comportamiento del sistema desde la perspectiva de todos
    los actores institucionales definidos en el \acuerdo; usar relaciones
    \textit{include/extend} apropiadamente
  & El diagrama cubre los tres subsistemas; los actores corresponden
    a los definidos en el \acuerdo; las relaciones \textit{include/extend}
    son semánticamente correctas; la sintaxis \plantuml\ es válida \\
  \midrule
  Especificaciones textuales de 3 casos de uso prioritarios
  (actor principal, precondición, flujo principal, flujos alternativos,
  postcondición; formato tabla)
  & Especificar de manera precisa el comportamiento esperado del sistema
    para los flujos de mayor complejidad
  & Las especificaciones son completas; los flujos alternativos cubren
    los escenarios de error más probables; la especificación es
    verificable y coherente con los -RF- del Sprint~2 \\
  \midrule
  Diagramas de secuencia en \plantuml\ para «Matricular Estudiante»
  y «Evaluar con Rúbrica» (con controlador, objetos de dominio y
  repositorios)
  & Modelar la interacción dinámica entre objetos del dominio;
    aplicar el patrón Controlador para separar la lógica de orquestación
    de las clases de dominio
  & Los diagramas reflejan fielmente el flujo del proceso; los mensajes
    corresponden a métodos reales de las clases del diagrama de dominio
    del Sprint~4; la separación controlador/dominio es clara \\
  \midrule
  Diagramas de estado para «Ciclo de vida del Estudiante» y «Ciclo
  de vida del Programa Académico» en \plantuml\
  & Modelar los estados y transiciones de las entidades del dominio
    con mayor variación de estado; identificar los invariantes de
    transición que el sistema debe hacer cumplir
  & Los diagramas son coherentes con el \acuerdo\ y el Decreto~1330
    de 2019; los estados terminales son correctos; las restricciones
    de transición están identificadas y justificadas \\
\end{longtable}

\subsection*{Actividad de reflexión: software generativo y diagramas de secuencia}
\label{act:ia-sprint5}

\begin{tcolorbox}[actividadIA,
  title={Actividad de reflexión (Sprint~5) —
         Software generativo y diagramas de secuencia}]

\textbf{Instrucción general.}
El estudiante construye el diagrama de secuencia para el caso de uso
«Matricular Estudiante» de manera individual, basándose en el modelo
de clases del Sprint~4 y los requisitos del Sprint~2. A continuación,
solicita al asistente generativo que genere el mismo diagrama, proporcionando
como contexto la especificación del caso de uso y las clases del dominio.
Se evalúa la coherencia de los mensajes generados con el modelo de clases.

\bigskip
\textbf{Preguntas orientadoras:}
\begin{enumerate}[label=\arabic*.]
  \item ¿Los mensajes del diagrama generado por el asistente corresponden
        a métodos reales de las clases del modelo de dominio del Sprint~4?
        Identifique al menos dos mensajes inventados o incoherentes con
        el modelo.
  \item ¿El asistente incorpora un controlador u objeto de servicio para
        orquestar la secuencia, o modela la lógica de orquestación dentro
        de las clases de dominio? ¿Cuál de las dos opciones es más coherente
        con el principio de responsabilidad única?
  \item ¿El asistente modela el flujo alternativo en que no hay cupos
        disponibles? Si no lo hace, ¿cómo reformularía el \textit{prompt}
        para obtener el diagrama completo?
  \item ¿El diagrama generado por el asistente es sintácticamente correcto
        en \plantuml\ y compila sin errores?
\end{enumerate}

\bigskip
\textbf{Rúbrica breve:}
\begin{itemize}
  \item \textbf{Producto} (60\%): coherencia del diagrama propio con el
        modelo de clases del Sprint~4, corrección de la secuencia de
        mensajes, inclusión del flujo alternativo, compilación exitosa
        del código \plantuml.
  \item \textbf{Proceso} (40\%): profundidad del análisis crítico del
        diagrama del asistente, identificación argumentada de
        inconsistencias conceptuales, reflexión sobre el valor del
        proceso de diseño manual frente al generado automáticamente.
\end{itemize}

\end{tcolorbox}
